\documentclass[11pt]{article}

\usepackage[letterpaper,margin=1in]{geometry}
\usepackage{amsmath,amssymb,amsthm,mathtools}
\usepackage{booktabs,array,enumitem,microtype,xurl}
\usepackage[hidelinks]{hyperref}

\hypersetup{
  pdftitle={Exact repair obstructions around a 64-modular Hadamard matrix of order 668},
  pdfauthor={Alec Kriebel, with substantial assistance from ChatGPT 5.6 Sol},
  pdfsubject={Exact local obstructions and certified reductions for Eliahou's order-668 modular Hadamard seed},
  pdfkeywords={Hadamard matrix, base sequence, Golay quadruple, Turyn sequence, SAT certificate}
}

\allowdisplaybreaks
\setlist{itemsep=2pt,topsep=4pt}

\newtheorem{theorem}{Theorem}[section]
\newtheorem{proposition}[theorem]{Proposition}
\newtheorem{lemma}[theorem]{Lemma}
\newtheorem{corollary}[theorem]{Corollary}
\newtheorem{computationaltheorem}[theorem]{Computational Theorem}
\newtheorem{computationalproposition}[theorem]{Computational Proposition}
\theoremstyle{remark}
\newtheorem{remark}[theorem]{Remark}

\newcommand{\Z}{\mathbb Z}
\newcommand{\F}{\mathbb F}
\newcommand{\BS}{\operatorname{BS}}
\newcommand{\TU}{\operatorname{TU}}
\newcommand{\supp}{\operatorname{supp}}
\newcommand{\wt}{\operatorname{wt}}
\newcommand{\dist}{\operatorname{dist}}
\newcommand{\rev}{\operatorname{rev}}
\newcommand{\Norm}{\operatorname{N}}
\newcommand{\Cor}{\operatorname{Cor}}
\newcommand{\PAF}{\operatorname{PAF}}
\newcommand{\sha}{\textsc{sha-256}}

\title{Exact repair obstructions around a\\
64-modular Hadamard matrix of order 668}
\author{Alec Kriebel\\[2pt]\large with substantial assistance from ChatGPT 5.6 Sol}
\date{Research draft, 25 July 2026\\Not peer reviewed}

\begin{document}
\maketitle

\begin{abstract}
Eliahou constructed a $64$-modular Hadamard matrix of order $668$ from a
special binary quadruple of length $167$.  An exact repair within that
special form is equivalent to a base sequence in $\BS(84,83)$.  We give
several exact obstructions and reductions around the published seed.

First, if Eliahou's sequence $q=(+^{83},-^2,+^{81},-)$ is fixed, no choice
of the other sign sequence gives an exact quadruple.  A parity telescope
and even--odd decimation would otherwise produce an element of
$\TU(41)$; in addition to the known classification, we independently
exhaust $461$ normalized shards, comprising $57{,}543{,}021$ deterministic
search nodes and no solution.

Second, reduction of the four base rows modulo $z^{42}-1$ proves that any
exact repair changes at least $80$ of the $334$ base signs, and at least
$41$ signs in the natural $(s,q)$ coordinates.  At special distance $41$,
the only remaining form is $39$ changes in $s$ and one reciprocal pair of
changes in $q$.  Root filters leave $21$ long and $18$ short reciprocal
pairs.  Reduction modulo $z^{42}+1$ removes endpoint orientation and
collapses these to $30$ canonical support problems.  A checked DRAT proof
excludes the boundary long case, and an exhaustive exact-integer census of
$3{,}710{,}853{,}316{,}608$ modular join rows excludes all nine short
cases.

The remaining twenty long cases are not excluded.  For them we prove an
exact endpoint-orientation cascade: the positive fold modulo $8$ is
redundant on the full anti-fold characteristic-two code, the modulo-$16$
layer is a fixed binary linear system, exact root conditions are four
cardinality equations, and the modulo-$32$ layer is quadratic.  We also
correct the scope of the two-fold method: the positive and negative
$42$-folds imply cyclic length-$84$ complementarity, not the original
aperiodic equations.  Forty-one causal equations restore exact
equivalence.  No $\BS(84,83)$, Golay quadruple, or Hadamard matrix of order
$668$ is constructed or ruled out.
\end{abstract}

\noindent\textbf{Verification and authorship disclaimer.}
This is an AI-assisted, unreviewed research artifact.  The mathematical
exploration, program development, certificate analysis, and manuscript
drafting used substantial assistance from OpenAI language models operating
through ChatGPT and Codex.  The named human author directed the project and
is responsible for any release, but describes himself as a non-specialist
and cannot independently certify every derivation or implementation.
The repository includes deterministic checkers and independently replayable
certificates where stated.  Those artifacts strengthen auditability; they
do not replace expert review, establish priority, or constitute a
proof-assistant formalization.

\section{Introduction}

A Hadamard matrix of order $n$ is a matrix
$H\in\{\pm1\}^{n\times n}$ such that
\[
 HH^{\mathsf T}=nI_n.
\]
The Hadamard conjecture predicts such a matrix whenever $4\mid n$.
Order $668$ is a prominent unresolved instance in current construction
tables~\cite{cati-pasechnik,sage}.  Eliahou constructed a
$64$-modular matrix at that order, satisfying
\[
 HH^{\mathsf T}\equiv 668I_{668}\pmod {64},
\]
from a structured length-$167$ quadruple and the Goethals--Seidel
array~\cite{eliahou-2025}.  Its failure of exact
complementarity is concentrated at thirteen aperiodic lags.
Eliahou's 2026 update is cited for bibliographic context only:
its metadata and abstract, but not its full text, were inspected during
this project~\cite{eliahou-2026}.  No claim below depends on that article.

The most direct local question is whether signs in that structured seed can
be changed to obtain exact complementarity.  Ramos, Hulak, and de Queiroz
have already recorded the basic translation of this question to
$\BS(84,83)$ and a raw base-row distance lower bound of $64$ in the
artifact release accompanying their work on multiplier
obstructions~\cite{ramos-2026,ramos-artifacts}.  The present paper starts
from the same standard base-sequence translation.  Its main additions are:
\begin{enumerate}
  \item a fixed-$q$ obstruction through $\TU(41)$, with an independent
        modern exhaustive certificate for the known endpoint;
  \item an adjacent-$42$ fold giving the stronger base-row distance bound
        $80$ and the exact special-coordinate boundary $41$;
  \item an orientation-free anti-fold, one checked DRAT exclusion, and a
        complete exact census of all nine short boundary cases; and
  \item an exact $2$-adic orientation cascade for the twenty open long
        cases, together with the necessary aperiodic scope correction.
\end{enumerate}

These conclusions are deliberately local.  The exact fixed-$q$ family is
closed, and ten of the thirty minimum-boundary anti-fold support cases are
closed.  Twenty long support cases remain open.  Nothing here implies
nonexistence of $\BS(84,83)$, of a nonspecial Golay quadruple of length
$167$, or of a Hadamard matrix of order $668$.

\section{Sequences and Eliahou's special form}

\subsection{Aperiodic norms and base sequences}

For a real sequence $X=(x_0,\ldots,x_{r-1})$, put
\[
 N_X(k)=\sum_{i=0}^{r-k-1}x_ix_{i+k}
 \qquad(0\leq k<r),
\]
and set $N_X(k)=0$ when $k\geq r$.  In polynomial notation,
$X(z)=\sum_i x_i z^i$, the Laurent polynomial
$X(z)X(z^{-1})$ has coefficient $N_X(|k|)$ at $z^k$.

A binary Golay quadruple of length $r$ is a quadruple
$(X_1,X_2,X_3,X_4)$ of sign sequences satisfying
\[
 \sum_{j=1}^4 N_{X_j}(k)=0
 \qquad(1\leq k<r).
\]
A base sequence in $\BS(m,n)$ is a quadruple $(A;B;C;D)$ of binary
sequences with
\[
 |A|=|B|=m,\qquad |C|=|D|=n,
\]
and
\[
 R_k:=N_A(k)+N_B(k)+N_C(k)+N_D(k)=0
 \qquad(1\leq k<m).
\]
We use the standard base-sequence notation and equivalences of
Đoković~\cite{djokovic-bs}.  Its zero-lag energy is $R_0=2m+2n$.
We use coefficientwise products of sign sequences without a separate
symbol.

\subsection{The special quadruple and its base rows}

Let $h\in\{\pm1\}^{167}$ be $+1$ on coordinates $0,\ldots,83$ and
$-1$ on coordinates $84,\ldots,166$.  Eliahou's special quadruple has
the form
\begin{equation}
 \mathcal G(s,q)=(s,hs,sq,hsq).                 \label{eq:special}
\end{equation}
The published sequence $q$ has alternating run lengths
\[
 q=(+^{83},-^2,+^{81},-).
\]
For a direct reconstruction of the seed, the run-length word of $s$ is
\begin{align*}
 &((4)^5,(2,1,1)^5,1,5,(4)^4,(2,1,1)^6,(4)^4,3,\\
 &\hspace{31mm}(1,2,1)^5,3,(4)^4,3,(1,2,1)^5),
\end{align*}
where the first run is positive and a parenthesized exponent denotes block
repetition.

Split at the sign change of $h$:
\begin{equation}
 A=s_{[0,84)},\quad B=(sq)_{[0,84)},\quad
 C=s_{[84,167)},\quad D=(sq)_{[84,167)}.       \label{eq:base-map}
\end{equation}

\begin{proposition}[Special-to-base-sequence translation]
\label{prop:translation}
The quadruple $\mathcal G(s,q)$ is exactly Golay if and only if
$(A;B;C;D)\in\BS(84,83)$.  The correspondence is bijective, with inverse
\[
 s=A\Vert C,\qquad q=(AB)\Vert(CD).
\]
\end{proposition}

\begin{proof}
At lag $k$, the summed correlation in \eqref{eq:special} is
\[
 \sum_i s_is_{i+k}
 (1+h_ih_{i+k})(1+q_iq_{i+k}).
\]
Terms crossing the $84+83$ split vanish.  The remaining value is twice
$R_k$ for the four rows in \eqref{eq:base-map}, and it vanishes
automatically for $k\geq84$.  The inverse follows coordinatewise.
\end{proof}

For the published base rows, the only nonzero positive-lag values of $R_k$
are
\[
\begin{array}{c|rrrrrrrrrrrrr}
k&4&8&12&16&26&30&34&38&42&46&50&54&58\\ \hline
R_k&-256&192&-128&64&-32&64&-96&128&-160&128&-96&64&-32.
\end{array}                                                \tag{2.1}
\]
The corresponding correlations of the four length-$167$ rows are twice
these values.  In particular, divisibility by $64$ does not imply exact
complementarity.  In the associated Goethals--Seidel matrix, Eliahou
records that every row is already orthogonal to $641$ of the other $667$
rows, with $26$ nonzero off-diagonal Gram entries~\cite{eliahou-2025}.

\section{The fixed-\texorpdfstring{$q$}{q} obstruction}
\label{sec:fixedq}

This section keeps the published $q$ fixed and permits arbitrary $s$.
The proof reduces the exact equations to a Turyn sequence of parameter
$41$, then uses both the known classification and an independent
enumeration.

\subsection{A parity telescope}

Write
\[
 X=(s_0,\ldots,s_{82}),\quad u=s_{84},\quad
 Y=(s_{85},\ldots,s_{165}),\quad v=s_{166}.
\]
The coordinate $s_{83}$ is isolated by the four runs of $q$.  Writing
$X=(x_0,\ldots,x_{82})$ and $Y=(y_0,\ldots,y_{80})$, exact
complementarity is equivalent to
\begin{align}
 N_X(k)+N_Y(k)&=0 &&(1\leq k\leq80),       \label{eq:fixed-main}\\
 x_0x_{81}+x_1x_{82}&=0,                   \label{eq:fixed-81}\\
 x_0x_{82}+uv&=0.                          \label{eq:fixed-82}
\end{align}

\begin{lemma}[Parity telescope]\label{lem:telescope}
Every solution of \eqref{eq:fixed-main}--\eqref{eq:fixed-82} satisfies
\begin{align}
 x_{82-i}&=(-1)^{i+1}x_i &&(0\leq i\leq82), \label{eq:Xreverse}\\
 y_{80-i}&=(-1)^iy_i     &&(0\leq i\leq80), \label{eq:Yreverse}\\
 u&=v.                                             \label{eq:uv}
\end{align}
\end{lemma}

\begin{proof}
For $0\leq k\leq81$, let $P_k$ be the product of all sign terms in
$N_X(k)+N_Y(k)$, with the $Y$ product empty at $k=81$.  At every positive
lag the zero sum contains $2(82-k)$ signs and exactly half are negative.
Hence
\[
 P_k=(-1)^k,\qquad 0\leq k\leq81.
\]
Cancellation in $P_{i+1}/P_i$ gives
\[
 (x_ix_{82-i})(y_iy_{80-i})=-1
 \qquad(0\leq i\leq80).
\]
Set $r_i=x_ix_{82-i}$ and $t_i=y_iy_{80-i}$.  Reflection gives
$r_i=r_{82-i}$ and $t_i=t_{80-i}$.  Comparing the equations at $i$ and
$80-i$ gives $r_i=r_{i+2}$.  The central products are
$r_{41}=t_{40}=1$, and $r_{40}t_{40}=-1$.  These values determine the two
parity classes of every $r_i$ and then every $t_i$, yielding
\eqref{eq:Xreverse} and \eqref{eq:Yreverse}.  Finally
$x_0x_{82}=-1$, so \eqref{eq:fixed-82} gives $uv=1$ and therefore $u=v$.
\end{proof}

\subsection{Even--odd decimation}

Define
\begin{align*}
 E&=(x_0,x_2,\ldots,x_{82}), & |E|&=42,& \rev(E)&=-E,\\
 O&=(x_1,x_3,\ldots,x_{81}), & |O|&=41,& \rev(O)&= O,\\
 P&=(y_0,y_2,\ldots,y_{80}), & |P|&=41,& \rev(P)&= P,\\
 Q&=(y_1,y_3,\ldots,y_{79}), & |Q|&=40,& \rev(Q)&=-Q,
\end{align*}
and let $\widehat Q=(u,Q,v)$, of length $42$.  Since $Q$ is skew and
$u=v$, the two endpoint contributions to $N_{\widehat Q}(k)$ cancel for
$1\leq k\leq40$.  The even-lag equations above, including the endpoint
equation, become
\[
 N_E(k)+N_{\widehat Q}(k)+N_O(k)+N_P(k)=0
 \qquad(1\leq k\leq41).
\]
Thus
\[
 (E;\widehat Q;O;P)\in\BS(42,41),
\]
with the first long row skew and the first short row symmetric.  Under the
standard definition this is an element of $\TU(41)$
~\cite{djokovic-quadruples,edmondson}.

\subsection{Independent exhaustion of \texorpdfstring{$\TU(41)$}{TU(41)}}

The nonexistence of $\TU(41)$ is already contained in the exhaustive
classification of Edmondson, Seberry, and Anderson~\cite{edmondson}.
For an independently reproducible endpoint, we also performed a fresh
enumeration.

The standard normalized form at odd parameter $41=2\cdot20+1$ is
\begin{equation}
\label{eq:normalized-tu}
\begin{aligned}
 A&=(1,1,a_2,\ldots,a_{20},-a_{20},\ldots,-a_2,-1,-1),\\
 B&=(1,1,b_2,\ldots,b_{20},-b_{20},\ldots,-b_2,-1, 1),\\
 C&=(1,c_1,\ldots,c_{19},c_{20},c_{19},\ldots,c_1,1),\\
 D&=(1,d_1,\ldots,d_{19},d_{20},d_{19},\ldots,d_1,1).
\end{aligned}
\end{equation}
This normalized search covers the object produced by the decimation above.
Definition 7(i) of Edmondson--Seberry--Anderson states that every odd
Turyn sequence---a base sequence whose first long row is skew and whose
first short row is symmetric---has a representative of
\eqref{eq:normalized-tu} under the standard sign, reversal, and
equal-length interchange normalizations, all of which preserve the summed
aperiodic norm~\cite{edmondson}.  The decimated quadruple
$(E;\widehat Q;O;P)$ has exactly those skew/symmetric types.  Hence a
hypothetical decimated solution would supply a normalized representative
enumerated below.  This bridge uses the cited normalization theorem; the
certificate independently exhausts its normalized representatives rather
than formalizing the equivalence theorem itself.

There are $78$ free signs.  The interchange symmetry of $C,D$ is removed
by lexicographic order.  Evaluation at $z=1$ gives
\[
 \sum(A)^2+\sum(B)^2+\sum(C)^2+\sum(D)^2=166.
\]
Here $\sum(A)=0$ and $\sum(B)=2$, so
\[
 |\sum(C)|=|\sum(D)|=9.
\]

\begin{computationaltheorem}[Independent $\TU(41)$ certificate]
\label{thm:tu41}
No normalized quadruple in \eqref{eq:normalized-tu} satisfies all $41$ positive-lag
base-sequence equations.  The deterministic outside-in enumeration has
\[
 461\ \text{complete shards},\qquad
 57{,}543{,}021\ \text{search nodes},\qquad
 0\ \text{solutions}.
\]
\end{computationaltheorem}

\paragraph{Certificate method.}
After symbolic cancellation, the high-lag equations introduce the
variables in small outside-in blocks.  A branch is retained only if every
newly complete high-lag equation is zero.  For a lower lag with fixed
contribution $F$ and $U$ unresolved product terms, a branch is rejected
only if $|F|>U$ or $F+U$ is odd.  These are necessary bounds even when the
unresolved products are dependent.  After five recursion steps, $19$
signs have been assigned and exactly $461$ prefixes survive.  An
independent Python program exhausts all $2^{19}$ prefix assignments and
verifies this cover.  Each shard then evaluates all $41$ equations at
every leaf.  Small cases $n=3,7,9$ reproduce the known pattern
present, present, absent.  The source, cube cover, shard index, and summary
are joined by hashes in the certificate.

\begin{theorem}[Fixed-$q$ obstruction]\label{thm:fixedq}
Keep Eliahou's published sequence
\[
 q=(+^{83},-^2,+^{81},-)
\]
and the special form $\mathcal G(s,q)$.  No binary sequence $s$ makes
$\mathcal G(s,q)$ exactly Golay.
\end{theorem}

\begin{proof}
A solution would yield an element of $\TU(41)$ by
Lemma~\ref{lem:telescope} and the even--odd decimation.  This contradicts
either the published classification~\cite{edmondson} or the independent
certificate in Theorem~\ref{thm:tu41}.
\end{proof}

\begin{remark}
The obstruction is not a two-squares argument.  The relevant row-sum
identity is
\[
 \sum(C)^2+\sum(D)^2=162=9^2+9^2.
\]
The final step is the exhaustive nonexistence of $\TU(41)$.
\end{remark}

\section{The adjacent-\texorpdfstring{$42$}{42} fold and exact distance}
\label{sec:distance}

\subsection{A cyclic image of \texorpdfstring{$\BS(84,83)$}{BS(84,83)}}

For a row $X$ of length $84$ or $83$, fold modulo $42$:
\[
 F_X(j)=X_j+X_{j+42},\qquad 0\leq j<42,
\]
where the missing second endpoint in a short row contributes zero.
Let $P_k$ be the summed periodic correlation of the four folded rows.
Ordinary reduction of the four-row norm modulo $z^{42}-1$ gives
\begin{align}
 P_0&=R_0+2R_{42},\nonumber\\
 P_k&=R_k+R_{42-k}+R_{42+k}+R_{84-k}
       &&(1\leq k\leq20),\label{eq:plusfold}\\
 P_{21}&=2(R_{21}+R_{63}).\nonumber
\end{align}
Consequently every $\BS(84,83)$ has a periodically complementary
length-$42$ integer image of energy $334$.

For Eliahou's base rows the folds are
\begin{equation}
 F_A=-2z^{41},\qquad F_B=0,\qquad
 F_C=2+z^{41},\qquad F_D=-2z^{40}+z^{41}.        \label{eq:seed-folds}
\end{equation}
Their summed cyclic norm is the constant $14$.  In particular, all
nonconstant fold correlations already vanish: the thirteen residual
aperiodic lags cancel in four triples under \eqref{eq:plusfold}.

\subsection{The distance-\texorpdfstring{$80$}{80} theorem}

Across the four base rows there are
\[
 42+42+41+41=166
\]
disjoint separation-$42$ pairs and two unpaired short-row signs.  If $E$
of those pairs have equal endpoints, the folded energy is
\begin{equation}
 2+4E.                                               \label{eq:fold-energy}
\end{equation}
An exact target has energy $334$, so it has exactly $83$ equal pairs.
The seed in \eqref{eq:seed-folds} has only three: row $A$ at cell $41$,
row $C$ at cell $0$, and row $D$ at cell $40$.

\begin{theorem}[Base and special distance bounds]\label{thm:distance}
Let $\mathcal E$ be Eliahou's labelled base quadruple.
\begin{enumerate}
  \item Every $(A;B;C;D)\in\BS(84,83)$ differs from $\mathcal E$ in at
        least $80$ of its $334$ base signs.
  \item Every exact special pair $(s,q)$ differs from the published
        $(s_0,q_0)$ in at least $41$ of its $334$ special-coordinate
        signs.
\end{enumerate}
\end{theorem}

\begin{proof}
Changing one base sign toggles the equality status of at most one
separation-$42$ pair.  Raising the number of equal pairs from $3$ to $83$
therefore takes at least $80$ changes, proving the first assertion.

At one coordinate, changing $s$ alone changes both associated base signs;
changing $q$ alone or changing both $s$ and $q$ changes one base sign.
Thus special distance $d$ produces base distance at most $2d$, and the
first assertion gives $d\geq40$.  Equality $d=40$ would require all
$40$ changes to be $s$-only, leaving $q=q_0$.  This is impossible by
Theorem~\ref{thm:fixedq}, so $d\geq41$.
\end{proof}

Equality in the first part is rigid.  All three seed-equal pairs and the
two short singletons remain fixed; exactly $80$ of the $163$ seed-opposite
pairs become equal; and exactly one endpoint is changed in each selected
pair.  If $F_0$ denotes the four folds in \eqref{eq:seed-folds}, write
$F=F_0+2G$.  Then $G$ is a four-row ternary word supported on exactly
$80$ seed-opposite cells and satisfies the exact group-ring equation
\begin{equation}
 \sum_r\bigl(F_{0,r}G_r^*+G_rF_{0,r}^*\bigr)
 +2\sum_r G_rG_r^*=160
 \quad\text{in }\Z[C_{42}].                         \label{eq:min-shell}
\end{equation}
The sign of every nonzero coefficient of $G$ identifies which endpoint
was changed.  Equation~\eqref{eq:min-shell} is necessary and sufficient
for the adjacent-$42$ cyclic image on the minimum shell, but not for the
original aperiodic equations.

\section{The special-distance-\texorpdfstring{$41$}{41} boundary}
\label{sec:boundary}

\subsection{The reciprocal \texorpdfstring{$q$}{q} skeleton}

Write $q=g\Vert h$, with $|g|=84$ and $|h|=83$.  Exact
complementarity forces the number of active equal-$q$ pairs at every lag to
be even.  Row reduction over $\F_2$ gives the following closed form.

\begin{lemma}[Reciprocal skeleton]\label{lem:skeleton}
Every exact special quadruple has
\begin{align}
 g&=(g_0,g_1,\ldots,g_{41},g_{41},\ldots,g_1,-g_0),\nonumber\\
 h&=(h_0,h_1,\ldots,h_{40},h_{41},h_{40},\ldots,h_1,h_0).
                                                               \label{eq:skeleton}
\end{align}
Conversely, every $q$ of this form makes the active-pair count even at
every lag.
\end{lemma}

\begin{proof}
Use bits for the signs of $q$.  Modulo two, the indicator that two signs
are equal is $1+q_i+q_j$.  If $e_k$ is the parity equation at lag $k$,
then elementary row operations give
\[
\begin{array}{ll}
 e_{83}:&g_0+g_{83}=1,\\
 e_j+e_{83-j}:&g_j+g_{83-j}=0\quad(1\leq j\leq41),\\
 e_{j+1}+e_{83-j}:&h_j+h_{82-j}=0\quad(0\leq j\leq40).
\end{array}
\]
These $83$ independent equations are precisely \eqref{eq:skeleton}.
Reversing the row operations proves the converse.
\end{proof}

Relative to the published $q_0$, the skeleton has exactly one possible
weight-one change, at the center of the short block (global coordinate
$125$).  Every legal weight-two change is one reciprocal pair; there are
$42$ long and $41$ short pairs.

\subsection{Classification at equality}

At one special coordinate let
\[
 a=\#(s\text{-only changes}),\quad
 b=\#(q\text{-only changes}),\quad
 c=\#(\text{changes in both }s\text{ and }q).
\]
Then
\[
 d_{\rm special}=a+b+2c,\qquad
 d_{\rm base}=2a+b+c.                                \label{eq:costs}
\]

\begin{theorem}[Complete special-distance-$41$ split]
\label{thm:boundary}
An exact repair at special distance $41$, if one exists, has exactly
\[
 39\ s\text{-only changes}
 \quad+\quad
 2\ q\text{-only changes},
                                                        \label{eq:only-boundary}
\]
and the two $q$ changes form one reciprocal pair.

Before root filtering there are $80$ shell-compatible reciprocal pairs.
The evaluations at $z=1$ and $z=-1$ leave exactly $39$ pairs:
$21$ long and $18$ short.  Every surviving pair has exactly two compatible
joined root profiles.
\end{theorem}

\begin{proof}
Combining $d_{\rm special}=41$ with $d_{\rm base}\geq80$ in
\eqref{eq:costs} leaves only
\[
 (a,b,c)=(41,0,0),\ (40,1,0),\ (39,2,0).
\]
The first keeps $q$ fixed and is impossible by
Theorem~\ref{thm:fixedq}.  In the second, the unique $q$ change is the
short center.  If $\lambda$ and $\sigma$ denote the ordinary signed sums
of the common long and short fold changes, the root energy equation is
\[
 \lambda(\lambda+1)+\sigma^2=41,
\]
or
\[
 (2\lambda+1)^2+(2\sigma)^2=165.
\]
This has no integer solution because both $3$ and $11$, primes congruent
to $3$ modulo $4$, occur to odd exponent in $165$.  Hence only
\eqref{eq:only-boundary} remains, and Lemma~\ref{lem:skeleton} makes the
two $q$ changes reciprocal.

Three of the $83$ reciprocal pairs are incompatible with equality in the
base distance bound, leaving $80$.  For a reciprocal $q$ pair, let
$H(z)$ be its signed two-term contribution to the active product row after
the adjacent fold, and define
\[
 h_+=H(1),\qquad h_-=H(-1).
\]
The exact $z=\pm1$ calculation gives
the following signed two-term fold distributions:
\[
\begin{array}{c|rrrr}
\text{long }(h_+,h_-)&(-2,0)&(0,2)&(0,-2)&(2,0)\\
\text{count}&6&15&15&6\\ \hline
\text{short }(h_+,h_-)&(2,-2)&(0,0)&(-2,2)&\\
\text{count}&10&18&10&
\end{array}
\]
Only the first two long types and the middle short type have compatible
root energy and parity profiles.  Their counts are $6+15+18=39$.
Solving the two exact sum-of-squares systems gives two profiles per pair.
\end{proof}

Writing a profile as
$((L(1),S(1)),(L(-1),S(-1)))$, the two profiles are
\begin{equation}
\label{eq:root-profiles}
\begin{array}{c|cc}
\text{pair type}&\text{profile 0}&\text{profile 1}\\ \hline
\text{long }(-2,0)&((-3,4),(-5,-4))&((6,-5),(4,5))\\
\text{long }(0,2)&((-4,-5),(-6,5))&((5,4),(3,-4))\\
\text{short }(0,0)&((-4,-5),(4,5))&((5,4),(-5,-4)).
\end{array}
\end{equation}

\section{The orientation-free anti-fold}
\label{sec:antifold}

\subsection{Reduction modulo \texorpdfstring{$z^{42}+1$}{z42+1}}

For a base row define its anti-fold
\[
 D_X(j)=X_j-X_{j+42},\qquad 0\leq j<42,
\]
again treating the missing short endpoint as zero.  Let $Q_k$ be the
summed negacyclic correlation of the four anti-folds.  Reduction modulo
$z^{42}+1$ gives
\begin{align}
 Q_0&=R_0-2R_{42},\nonumber\\
 Q_k&=R_k-R_{42-k}-R_{42+k}+R_{84-k}
       &&(1\leq k\leq20),                       \label{eq:antifold}\\
 Q_{21}&=0,\qquad Q_{42-k}=-Q_k.\nonumber
\end{align}
Every $\BS(84,83)$ therefore satisfies
\[
 \sum_X D_X(z)D_X(z^{-1})=334
 \quad\text{in }\Z[z]/(z^{42}+1).                    \label{eq:anti-norm}
\]

For a seed-opposite pair $(x,-x)$, changing the lower endpoint gives
$(-x)-(-x)=0$, whereas changing the upper endpoint gives $x-x=0$.
Thus endpoint orientation disappears from the anti-fold.  At the boundary
in Theorem~\ref{thm:boundary}, the $39$ changes in $s$ select $39$ cells;
each selected long cell is zeroed in both long anti-fold rows, and each
selected short cell is zeroed in both short anti-fold rows.  The fixed
reciprocal $q$ pair zeroes two further entries in the active product row.
The anti-fold energy drops by exactly
\[
 80\cdot4=320,
\]
from the seed value $654$ to the required $334$.  The remaining task is
the set of $20$ noncentral equations in \eqref{eq:anti-norm}.

The $21$ long reciprocal pairs remain distinct.  The $18$ short pairs
occur in nine mirrored pairs,
\[
 (2,38),(4,36),\ldots,(18,22),
\]
and each mirrored pair gives the same orientation-free support problem.
There are therefore
\[
 21\text{ long cases}+9\text{ short cases}=30
\]
canonical anti-fold instances.  We label the long boundary instance
$L0$ as case $0$, the other long instances as cases $1,\ldots,20$, and
the short instances $S02,S04,\ldots,S18$ as cases $21,\ldots,29$.

\subsection{A checked DRAT exclusion}

Case $0=L0$ has $79$ available support variables; all other cases have
$78$.  In every case exactly $39$ support cells must be selected.

\begin{computationaltheorem}[Case-$0$ anti-fold exclusion]
\label{thm:case0}
The complete orientation-free anti-fold support problem for the long
reciprocal pair $L0$ is unsatisfiable.  This exclusion is independent of
the two adjacent-fold root profiles.
\end{computationaltheorem}

\paragraph{Encoding and certificate.}
The exact cardinality equation and the twenty integer anti-fold equations
were translated to CNF.  Every product of two retained-entry literals is
linked to an auxiliary variable by the three-clause Boolean AND
equivalence.  Signed pseudo-Boolean sums are encoded exactly.  Redundant
modulo-two and modulo-four Hensel chains are consequences of the integer
equations and therefore do not remove a genuine support.  The resulting
formula has
\[
 39{,}580\ \text{variables},\qquad
 127{,}589\ \text{clauses}.
\]
Its \sha\ digest is
\[
\texttt{f3eb29b1ea9c386e53b03726349fe0c38577d7e187b56aa19f86412c8749755d}.
\]
CaDiCaL 3.0.1 produced a binary DRAT proof in $200.17$ seconds.  The
$90{,}490{,}737$-byte compressed proof has digest
\[
\texttt{efd8abd9d80d50365822754f36345f368d7cff8f2740ca33b9cab7d5866aa519}.
\]
Independent \texttt{drat-trim} replay returned \texttt{VERIFIED}, using
$73{,}135{,}509$ resolution steps, $0$ RAT lemmas, $75.00$ seconds, and
$471.1$ MB peak resident memory.

The formal proof trace certifies the generated CNF.  The mathematical
interpretation additionally relies on the auditable encoder and on the
reduction to the support equations; neither is formalized in a proof
assistant.

\subsection{Complete census of the nine short cases}

The short cases have a common exact quotient structure.  Modulo two, the
$78$ support variables form $39$ pairs of equal syndrome columns.  The
affine quotient has rank $21$, leaving $2^{18}$ parity states.  For every
state, a phase-adjusted parity splits the characteristic-three interaction
graph into four cliques:
\[
 L_0,\ L_1,\ S_0,\ S_1.
\]
After conditioning the common central pair, sums from these four components
can be joined by exact packed modulo-six signatures.  Any survivor is
reconstructed as all $78$ support bits, checked against the exact
twenty-coordinate integer polynomial, and independently replayed in the
four physical anti-fold rows.

Spatial reflection acts freely because every quotient state has an odd
noncentral reflected pair.  The canonical gauge fixes the lowest-index odd
long-pair orientation.  Two exceptional quotient states, one in $S08$ and
one in $S14$, instead use the lowest odd short pair.  The reflected mate of
every retained representative is reconstructed, so the gauge loses no
support.

\begin{computationaltheorem}[All nine short cases]
\label{thm:short}
The complete exact anti-fold support problems for cases $21$ through $29$,
namely
\[
 S02,S04,S06,S08,S10,S12,S14,S16,S18,
\]
have no solution.
\end{computationaltheorem}

The frozen census is summarized below.
\[
\begin{array}{cclrrr}
\toprule
\text{case}&\text{label}&\text{join rows}&
\text{mod-6 survivors}&\multicolumn{2}{c}{\text{best residual}}\\
&&&&\#\ne0&\ell_1\\
\midrule
21&S02&412{,}316{,}860{,}416& 9{,}564{,}254&4&72\\
22&S04&412{,}316{,}860{,}416& 9{,}796{,}880&4&78\\
23&S06&412{,}316{,}860{,}416& 9{,}807{,}738&3&54\\
24&S08&412{,}317{,}646{,}848& 9{,}561{,}296&4&66\\
25&S10&412{,}316{,}860{,}416& 9{,}789{,}738&3&72\\
26&S12&412{,}316{,}860{,}416&10{,}533{,}216&4&66\\
27&S14&412{,}317{,}646{,}848&10{,}285{,}492&3&66\\
28&S16&412{,}316{,}860{,}416& 9{,}789{,}970&4&66\\
29&S18&412{,}316{,}860{,}416& 9{,}799{,}156&4&66\\
\midrule
\multicolumn{2}{r}{\text{total}}&
3{,}710{,}853{,}316{,}608&
88{,}927{,}740&&\\
\bottomrule
\end{array}                                           \tag{6.3}
\]
All $88{,}927{,}740$ modular survivors received both an exact
integer-polynomial check and an independent bit-packed physical replay.
The number of exact supports was zero.  The production was partitioned
into $2{,}304$ atomic ranges.  The completion certificate freezes every
case total, best witness, source and binary hash, the original
survivor-stream digest, and a timing-independent semantic range digest.
Without the underlying range manifests, those stream digests are
cryptographic commitments rather than independently inspectable survivor
lists.  The immutable companion release includes the production-manifest
archive, so the live verifier can check all $2{,}304$ records; an
independent full production replay remains the strongest check.

\begin{corollary}[Certified minimum-boundary frontier]
\label{cor:frontier}
Among the thirty canonical orientation-free anti-fold support problems at
special distance $41$, case $0$ and all nine short cases are impossible.
Only the twenty long cases $1,\ldots,20$ remain.
\end{corollary}

\begin{remark}
Corollary~\ref{cor:frontier} is not a distance-$41$ nonexistence theorem.
It closes ten support cases and leaves twenty.  It also says nothing about
repairs at larger special distance or about base sequences unrelated to
Eliahou's seed.
\end{remark}

\section{Exact orientation cascade for the open long cases}
\label{sec:orientation}

The twenty open long instances all have $78$ eligible support cells,
split into $39$ long and $39$ short cells.  Fix a support
$S$ of exact weight $39$ that satisfies the characteristic-two anti-fold
conditions and one of the two root profiles in \eqref{eq:root-profiles}.  The
remaining choice is which endpoint of each selected seed-opposite pair is
changed.

For a selected cell $i$, let $\delta_i$ be the change in the four positive
fold rows when the lower endpoint is chosen.  The upper endpoint contributes
$-\delta_i$.  Encode the upper choice by $e_i\in\F_2$.

\subsection{Modulo \texorpdfstring{$8$}{8} is redundant}

Let $C_k$ be the summed periodic correlation at cyclic lag $k$ of the
positive folds.  Exact expansion into constant, singleton, and pair
coefficients proves the following.

\begin{theorem}[Long-case orientation cascade]\label{thm:cascade}
For each open long case $1,\ldots,20$:
\begin{enumerate}
  \item For every support in the full $57$-dimensional affine anti-fold
        characteristic-two code, all positive-fold correlations vanish
        modulo $8$, independently of endpoint orientation.
  \item There is a support-independent binary matrix
        $M_c\in\F_2^{23\times78}$ such that a support $S$ selects its
        columns and the complete modulo-$16$ plus-fold and root-parity
        conditions are
        \[
        (M_c)_S e=b_{c,S,\rho}.                       \label{eq:mod16}
        \]
        The first $21$ rows encode $C_k/8\bmod2$ and the final two rows
        encode the long- and short-block root-orientation parities.
  \item On any consistent affine solution fiber of
        \eqref{eq:mod16}, the four exact root values are equivalent,
        after a fixed sign gauge, to four disjoint Hamming-weight
        equations.
  \item If a consistent fiber has free coordinates
        $u_1,\ldots,u_d$, the next $21$ plus-fold conditions modulo $32$
        are exact Boolean quadratics
        \[
        q_{c,S,\rho,k}(u)=0,\qquad 1\leq k\leq21.      \label{eq:mod32}
        \]
\end{enumerate}
\end{theorem}

\begin{proof}[Proof outline and exact checks]
For the first assertion, direct integer expansion proves that every
positive-fold correlation is divisible by $4$ for arbitrary support and
that $C_k/4\bmod2$ is affine in the $78$ support bits.  Restricting these
$21$ affine functions to a basis and particular point of each complete
$57$-dimensional anti-fold affine code makes every coefficient and constant
zero.  All $20$ cases and all $3{,}003$ second differences per case were
also reconstructed by an independent implementation.

Every orientation interaction coefficient
\[
 \langle\delta_i,T^k\delta_j\rangle+
 \langle\delta_j,T^k\delta_i\rangle
\]
is divisible by $8$.  After division by $8$ and reduction modulo $2$, the
orientation dependence is therefore affine, and the column belonging to
cell $i$ is independent of the other selected cells.  This constructs
$M_c$ and proves \eqref{eq:mod16}.

For a block/parity bin $G$, let $a_j\in\{\pm1\}$ be the signed root
contribution in the lower orientation and put
$n_j=e_j\mathbin{\mathsf{xor}}[a_j=-1]$.  A prescribed signed target $t_G$
is equivalent to
\[
 \sum_{j\in G}a_j(-1)^{e_j}
 =|G|-2\wt(n|_G)=t_G,
\]
which is exactly one cardinality equation.

Finally, substitute an affine parameterization of the solution space of
\eqref{eq:mod16}.  An orientation flip changes a positive-fold coefficient
by a multiple of $4$.  After division by $16$ and reduction modulo $2$,
all terms of degree at least three vanish.  Direct interpolation of every
constant, linear, and quadratic coefficient, followed by physical replay
at deterministic higher-weight points, proves \eqref{eq:mod32}.
\end{proof}

No universal lower bound on $\operatorname{rank}(M_c)_S$ is asserted.  In
the pinned audits, consistent supports usually had rank $22$ and hence
nullity $17$, with a small number of nullity-$18$ fibers.  Those counts are
finite diagnostics, not a theorem about all supports.

\begin{computationalproposition}[Sparse generation of the mod-$2$ support codes]
\label{prop:sparse-code}
For each of the twenty open long cases, the $57$-dimensional homogeneous
anti-fold characteristic-two code is generated by words of weight at most
$6$.  In cases $6$ and $14$, weight at most $4$ suffices.  After adjoining
the four block/parity profile checks, the homogeneous code has dimension
$55$ and is generated by words of weight at most $6$ in every case.
\end{computationalproposition}

\paragraph{Exact rank certificate.}
The verifier
\path{eliahou_long_orientation_cascade/audit_orientation_cascade.py}
constructs the complete binary check matrix: the anti-fold linear rows and
the all-ones weight-parity row, with the four block/parity rows adjoined in
the profile-conditioned calculation.  It encodes each coordinate column by
its syndrome.  Two unordered coordinate pairs with the same XOR syndrome
have a symmetric difference in the kernel of weight at most $4$.
The verifier groups all pairs by syndrome, inserts these differences into
an exact binary Gaussian basis, and records their rank.  If that rank is
not the full code dimension, it analogously groups all coordinate triples
by syndrome; differences of equal-syndrome triples have weight at most
$6$.  The resulting basis ranks are exactly $57$ in every unconditioned
case and $55$ in every conditioned case.  Pair differences already attain
rank $57$ in unconditioned cases $6$ and $14$, proving their stronger
weight-$4$ statement.  The ordered rank arrays and a semantic hash are
frozen in \path{EXPECTED_SAMPLE5_SUMMARY.json}.

This provides an exact connected local-move system, but it does not bound
the number of exact-weight-$39$ supports sharply enough for enumeration.

\subsection{Certified finite controls}

Two independent implementations used different deterministic sets of five
supports per case and profile.  Their exact finite results are:
\begin{equation}
\label{eq:bounded-controls}
\begin{array}{lrr}
\toprule
&\text{primary audit}&\text{red-team audit}\\
\midrule
\text{supports tested}&200&200\\
\text{mod-$16$ inconsistent}&5&5\\
\text{points in consistent fibers}&26{,}607{,}616&26{,}476{,}544\\
\text{points satisfying exact roots}&121{,}764&120{,}455\\
\text{mod-$32$ survivors}&22&24\\
\text{mod-$32$ survivors satisfying exact roots}&1&0\\
\bottomrule
\end{array}
\end{equation}
The one root-exact modulo-$32$ point in the primary audit was independently
replayed and failed modulo $64$.  Thus all $400$ pinned fixture evaluations
were excluded at the recorded stage.  Equation~\eqref{eq:bounded-controls} is an exact
statement about those fixed inputs only; it is not a sampling theorem,
frequency estimate, or exclusion of any long case.

\section{Why the two folds are not aperiodically complete}
\label{sec:scope}

The positive and negative folds at the same $42$ split do not reconstruct
the original aperiodic norm.  This point is essential to the scope of every
computational claim above.

\begin{proposition}[Two-fold scope correction]\label{prop:scope}
Let $R_k$ be the summed aperiodic correlations of four rows of lengths
$84,84,83,83$.  Exact positive-fold and anti-fold norm equations together
are equivalent to the cyclic length-$84$ conditions
\begin{equation}
 R_{42}=0,\qquad R_k+R_{84-k}=0
 \quad(1\leq k\leq41).                              \label{eq:cyclic84}
\end{equation}
They do not imply $R_k=0$ separately.

Adding the causal half
\begin{equation}
 R_{43}=R_{44}=\cdots=R_{83}=0                       \label{eq:causal}
\end{equation}
to \eqref{eq:cyclic84} is equivalent to all $83$ positive-lag aperiodic
equations.
\end{proposition}

\begin{proof}
Adding and subtracting the coefficient identities
\eqref{eq:plusfold} and \eqref{eq:antifold} yields exactly
\eqref{eq:cyclic84}.  The residual vector
\[
 R_1=2,\qquad R_{83}=-2,\qquad R_k=0\text{ otherwise}
\]
is a nonzero integral kernel element, so the implication to aperiodic
complementarity fails.  If \eqref{eq:causal} is added, then
\eqref{eq:cyclic84} successively forces $R_1,\ldots,R_{41}$ to vanish as
well, while $R_{42}=0$ is already present.
\end{proof}

There is a compact exact ternary formulation of each long boundary case.
For each eligible cell let
\[
 t_j\in\{-1,0,1\},\qquad u_j=t_j^2,
\]
where $t_j=0$ means no change and the two nonzero signs choose the endpoint
orientation.  Then the exact physical problem is
\begin{equation}
\begin{aligned}
 \sum_j u_j&=39,\\
 20\ \text{anti-fold equations}&\quad\text{in }u,\\
 21\ \text{positive-fold equations}&\quad\text{in }t,\\
 41\ \text{causal equations \eqref{eq:causal}}.&
\end{aligned}                                        \label{eq:ternary}
\end{equation}
Independent symbolic reconstruction verifies \eqref{eq:ternary} against
all original correlations in all twenty cases.  The direct aperiodic model
and the fold-plus-causal model each use exactly
\[
 5{,}928
 =1{,}482\,uu+1{,}482\,ut+1{,}482\,tu+1{,}482\,tt
\]
quadratic products.  The causal equations touch all of them.  Thus the
folds expose strong modular structure, but they do not secretly eliminate
the global aperiodic coupling.

\section{Reproducibility and trust boundary}
\label{sec:repro}

All paths below are relative to the public repository's
\texttt{hadamard\_668\_search} directory.  The accompanying paper README
gives exact commands and separates quick verification from complete
production replay.

\begin{center}
\scriptsize
\begin{tabular}{>{\raggedright\arraybackslash}p{0.21\textwidth}
                >{\raggedright\arraybackslash}p{0.28\textwidth}
                >{\raggedright\arraybackslash}p{0.35\textwidth}}
\toprule
\textbf{Claim} & \textbf{Primary artifact} & \textbf{Independent check}\\
\midrule
Fixed-$q$ reduction &
\path{FIXED_Q_OBSTRUCTION.md} &
\path{verify_fixed_q_obstruction.py}\\
Distance $80$, boundary $41$, root frontier &
\path{ELIAHOU_ADJACENT42_REPAIR.md} &
\path{verify_eliahou_adjacent42_repair.py}\\
Anti-fold reduction &
\path{ELIAHOU_ANTIFOLD42.md} &
\path{verify_eliahou_antifold42.py}\\
$\TU(41)$ exhaustion &
\path{tu41_certificate/certificate.json} &
\path{verify_manifest.py}, \path{verify_cube_cover.py},
\path{test_regressions.py}\\
Case-$0$ UNSAT &
\path{output/antifold42_q0_proof/certificate.json} &
\path{verify_eliahou_antifold_q0_proof.py}; full
\path{drat-trim} replay\\
Nine short cases &
\path{eliahou_short_block_census/NINE_CASE_COMPLETION_CERTIFICATE.json} &
\path{verify_nine_case_completion.py}; independent model and physical
replays\\
Long orientation cascade &
\path{eliahou_long_orientation_cascade/} &
separate \path{eliahou_long_orientation_redteam/} implementation\\
\bottomrule
\end{tabular}
\end{center}

The trust levels are not identical.
\begin{itemize}
  \item The algebraic reductions are short exact-integer identities with
        dependency-free reconstruction tests.
  \item The $\TU(41)$ result is an exhaustive C++ enumeration with an
        independently exhausted prefix cover.  It trusts the displayed
        normalization, the small programs, and the compiler.
  \item The case-$0$ formula has an independently checked DRAT proof.
        It additionally trusts the formula generator's faithful encoding of
        the mathematical support problem.
  \item The nine-case certificate freezes the full production semantics,
        hashes, counts, and survivor streams.  Bounded independent replays
        reconstruct quotient states and physical rows.  Repeating all
        $3.71$ trillion joins is possible with the resumable production
        driver but is not part of the quick verifier.
  \item The long-case statements labeled as theorems are complete symbolic
        audits over their stated affine spaces.  The $200$-support tables
        are explicitly bounded controls and have no extrapolated scope.
\end{itemize}

No solver timeout, interrupted search, or unproved solver
\texttt{INFEASIBLE} status is used as a theorem in this manuscript.  In
particular, an older uncertified local-radius computation is intentionally
omitted.

\section{Discussion}

The fixed-$q$ theorem closes a natural but very narrow repair family.
The adjacent fold then shows that ``nearby'' in the published special
coordinates is already a global condition: an exact target cannot occur
before distance $41$, and equality forces the rigid $39+2$ form of
Theorem~\ref{thm:boundary}.  The anti-fold converts that boundary into a
finite collection of orientation-free support equations.  Ten of the
thirty canonical instances are now exactly excluded.

The open part is not the endpoint-orientation layer by itself.
Within the $400$ bounded fixture evaluations in
\eqref{eq:bounded-controls}, Theorem~\ref{thm:cascade} reduced a consistent
fixed-support orientation problem from $2^{39}$ raw choices to a
modulo-$16$ affine fiber usually of size about $2^{17}$ before the exact
root and quadratic modulo-$32$ filters.  No $2^{17}$ fiber-size claim is
made for untested supports or for a complete long case.  The
bottleneck is the enormous collection of admissible supports and the
aperiodic causal coupling in Proposition~\ref{prop:scope}.  A successful
continuation appears to require a support-level theorem or a method that
couples support and orientation algebraically; a literal support-first
census remains out of reach on ordinary hardware.

The known empty status of $\TU(41)$ is not claimed as new.  The independent
enumeration is valuable as a modern, inspectable reproduction.  The basic
special-form/$\BS(84,83)$ equivalence and a weaker distance bound were
already present in the companion release to
Ramos--Hulak--de Queiroz~\cite{ramos-artifacts}.  We did not locate the
distance-$80$ fold, special-distance-$41$ classification, anti-fold
certificates, or long orientation cascade in the cited public sources.
That literature search is not a guarantee of priority.

\section*{Code and data availability}

The immutable companion release is
\begin{center}
\url{https://github.com/AlecKriebel/Math/releases/tag/h668-research-checkpoint-v1.0.0}.
\end{center}
The paper package identifies the exact source paths and hashes used for
this draft.  The companion release includes the case-$0$ compressed DRAT
artifact and a production-manifest archive for the nine short cases.
Large artifacts are additionally identified by cryptographic digests.
No external communication or independent peer validation occurred during
preparation of this draft.

\section*{AI assistance and responsibility}

The research program was conducted by Alec Kriebel with extensive
assistance from ChatGPT 5.6 Sol, operating through ChatGPT and Codex.
AI assistance included proposing
reductions, deriving and checking algebra, writing and optimizing exact
enumerators, constructing certificate verifiers, auditing scope, and
drafting this manuscript.  The human author selected the research
directions and requested the explicit scope and verification disclosures.
Language-model output can contain subtle mathematical or implementation
errors.  Readers should therefore treat every claim as provisional until
the proofs, encodings, and computations receive independent expert review.
The software and certificates are provided to make that review feasible,
not to transfer responsibility to the tools.

\begin{thebibliography}{99}
\small
\setlength{\itemsep}{0.35em}

\bibitem{cati-pasechnik}
M. Cati and D. V. Pasechnik,
\emph{A database of constructions of Hadamard matrices},
arXiv:2411.18897 (2024);
\url{https://arxiv.org/abs/2411.18897}.

\bibitem{djokovic-quadruples}
D. Ž. Đoković,
Aperiodic complementary quadruples of binary sequences,
\emph{J. Combin. Math. Combin. Comput.} \textbf{27} (1998), 3--31;
\url{https://combinatorialpress.com/article/jcmcc/Volume%20027/vol-27-paper%201.pdf}.

\bibitem{djokovic-bs}
D. Ž. Đoković,
Classification of base sequences $\BS(n+1,n)$,
\emph{Int. J. Combin.} (2010), Article ID 851857;
\url{https://arxiv.org/abs/1002.1414}.

\bibitem{edmondson}
G. M. Edmondson, J. Seberry, and M. R. Anderson,
On the existence of Turyn sequences of length less than 43,
\emph{Math. Comp.} \textbf{62} (1994), 351--362;
\url{https://doi.org/10.1090/S0025-5718-1994-1203733-8}.

\bibitem{eliahou-2025}
S. Eliahou,
A 64-modular Hadamard matrix of order 668,
\emph{Australas. J. Combin.} \textbf{93}(2) (2025), 422--427;
\url{https://ajc.maths.uq.edu.au/pdf/93/ajc_v93_p422.pdf}.

\bibitem{eliahou-2026}
S. Eliahou,
An update on modular Hadamard matrices,
\emph{J. Algebraic Combin.} \textbf{64} (2026), Article 1;
\url{https://doi.org/10.1007/s10801-026-01544-5}.
Only its bibliographic metadata and abstract were inspected for this draft;
the full text was not inspected.

\bibitem{ramos-2026}
A. F. Ramos, D. B. Hulak, and R. J. G. B. de Queiroz,
\emph{Multiplier obstructions for Legendre pairs of length 333},
arXiv:2607.20765 (2026);
\url{https://arxiv.org/abs/2607.20765}.

\bibitem{ramos-artifacts}
A. F. Ramos, D. B. Hulak, and R. J. G. B. de Queiroz,
\emph{Hadamard-668 multiplier-obstruction proof artifacts},
companion repository and mod-$64$ report (2026);
\url{https://github.com/Arthur742Ramos/hadamard-668-multiplier-obstructions}.

\bibitem{sage}
SageMath,
\emph{Hadamard matrix constructions and current unknown orders};
\url{https://doc.sagemath.org/html/en/reference/combinat/sage/combinat/matrices/hadamard_matrix.html}.

\end{thebibliography}

\end{document}
